\input{preamble}
\begin{document}

\title{Geometric Stability Conditions and
Group Actions}
\maketitle

\noindent
Minicourse by Hannah Dell (University of
Bonn, Germany), CIMPA school Florian\'opolis
2025.

\medskip\noindent
Hannah's notes at
\href{https://www.hannahdell.me/cimpa2025}
{hannahdell.me/cimpa2025}

\medskip\noindent
These notes at
\href{%
http://github.com/danimalabares/cimpa-floripa
}{%
github.com/danimalabares/cimpa-floripa
}

\bigskip\noindent

{\bf Abstract.} Group actions on categories
arise naturally from symmetries of varieties
and quivers, but how does this interact with
Bridgeland stability? In the frist half of
this course we will introduce equivariant
categories, which generalises the category of
equivariant sheaves. Then we will show there
is a correspondence between stability
conditions on a category with a fniite group
action, and stability conditions on the
equivariant category -- this will also play a
role in Xiaolei Zhao's course. We will use
this to produce stability conditions on
quotient varieties (and stacks).

In the second half of the course, we will
apply this to study open questions about the
geometry of the stability manifold: in
particular, we will discuss "geometric
stability conditions" -- those for which all
skyscraper sheaves of points are stable. In
practice, these are constructed using slope
stability for sheaves. Some varieties have
only geometric stability conditions, whereas
in other cases, there are more (for example
if there is an equivalence with quiver
representations). Lie Fu, Chunyi Li, and
Xiaolei Zhao were the frist to provide a
general result explaining this phenomenon. In
particular, they showed that if a variety has
a fniite map to an abelian variety, then all
stability conditions are geometric. We will
test the converse on free quotients of
abelian varieties by fniite groups, including
Beauville-type and bielliptic surfaces. This
is based on joint work with Edmund Heng and
Anthony Licata.

\bigskip\noindent
\tableofcontents
\bigskip\noindent

\section*{Motivation}
\label{section-motivation}

The Bridgeland stability machine:

$$
[diagram]
$$
\noindent
{\bf Question 1.} How does the geometry of
$\text{Stab}(X)$ relate to $X$?
\noindent
{\bf Question 2.} How to invariants (of $X$
or $\ddot\cup $) behave under (finite) group
actions?

\medskip\noindent
{\bf Goal of these lectures.} Answer both by
studying ``free quotients'' (and see lots of
examples along the way!)

\section{Geometric stability conditions on
surfaces}
\label{gscag-label-001}
\noindent
{\bf Setup.} $X$ smooth projective surface
over $\mathbb{C}$.
$$
\begin{aligned}
\text{NS}(X)&:=\text{Pic}(X)/\text{Pic}^0(X)
\qquad
\text{Neron-Severi group}\\
\text{NS}_{\mathbb{R}}(X)& :=\text{NS}(X)
\otimes \mathbb{R}\\
\text{Amp}_{\mathbb{R}}(X)&:=\text{ ample
cone, i.e. $H \in
\text{NS}_{\mathbb{R}}(X)$ such that $H^2>0$}
\end{aligned}
$$

\begin{definition}
\label{definition-slope-and-stability}
Let $0 \neq  E \in \text{Coh}(X)$ and $H \in
\text{Amp}_{\mathbb{R}}(X)$,
\begin{enumerate}
\item The {\it $H$-slope} of $E$ is
$$
\mu_H(E):=
\begin{cases}
+\infty\qquad &\text{ch}_0(E)=0 \\
\frac{H\cdot \text{ch}_1(E)}{H^2
\text{ch}_0(E)}\qquad &\text{otherwise}
\end{cases}
$$
\item $E$ is {\it $H$-(semi)stable} if $0
\neq  F \subset E$ implies
$$
\mu_H(F) < \mu_H(E/F)\qquad (\mu_H\leq
\mu_H(E/F)
$$
\end{enumerate}
\end{definition}

\begin{example}
\label{example-semistable-condition}

\end{example}

\noindent
{\bf Problem.} [missing]

\begin{definition}
\label{gscag-label-002}
A path $\sigma=(\mathcal{A},Z)$ is a {\it
(numerical) Bridgeland stability condition}
for $D^b(X)$ if
\begin{enumerate}
\item $\mathcal{A} \subset D^b(X)$ is the
heart of an bounded $t$-structure.
\item $Z:K_0(X) \to \mathbb{C}$ homomorphism
(the {\it central charge}) such
that
\begin{itemize}
\item $0 \neq E \in \mathcal{A} \implies
Z(E) \in \mathbb{H}$, where
$\mathbb{H}$ is the upper half plane in
$\mathbb{C}$.
\item HN filtrations exist.
\item Factors via $\text{Knum}(X)$.
\end{itemize}
\item Support property.
\end{enumerate}
\end{definition}

Let $X$ be a surface, and fix  $H \in
\text{Amp}_\mathbb{R}(X)$ and $\beta \in
\mathbb{R}$. Define
$$
\begin{aligned}
\tau_{M,\beta}&:=\{E \in \text{Coh}(X):E
\overset{\text{surj.}}{\to} Q \neq 0
\implies \mu_M(Q)>\beta\}\\
\mathcal{F}_{H,\beta}&:=\{E \in
\text{Coh}(X):0 \neq  F \subset E
\implies \mu_M(F) \leq \beta\}
\end{aligned}
$$

\begin{exercise}
\label{exercise-torsion-pair}
$(\tau_{M,\beta},\mathcal{F}_{H,\beta})$ is a
torsion pair in $\text{Coh}(X)$.
\end{exercise}

\begin{definition}
\label{definition-Coh-H-beta}
$$
\text{Coh}^{H,\beta}(X):=\left\{ E \in
D^b(X):
\substack{H^0(E) \in \tau_{H,\beta} \\
H^{-1}(E) \in \mathcal{F}_{H,\beta}\\
H^i(E)=0} \right\}
$$
\end{definition}

\begin{exercise}
\label{exercise-new-slope}

\end{exercise}

There are notes available about the
following:

\begin{theorem}[Bridgeland '08,
Arcara-Bertam, Macri-Schmidt]
\label{theorem-Bridgeland}
For $\alpha\gg 0$,
$$
\mathbb{C}\cdot \{\sigma_{H,D,\alpha,\beta}=
(\text{Coh}^{H,B}(X),Z_{H,D,\alpha,\beta})\}
$$
is a continuous family of Bridgeland
stability conditions.
\end{theorem}

\noindent
{\bf Question 3.} Is this all of
$\text{Stab}(X)$?

\begin{exercise}
\label{exercise-O-is-stable}
$\mathcal{O}_x$ is
$\sigma_{H,D,\alpha,\beta}$-stable for all
$x\in X$.
\end{exercise}

\begin{definition}
\label{gscag-label-003}
$X$ smooth projective variety, $\sigma \in
\text{Stab}(X)$ is {\it geometric} if for all
$x \in X$, $\mathcal{O}_x$ is
$\sigma$-stable. Write
$\text{Stab}^{\text{geo}}(X)=$ all geometric
stability conditions.
\end{definition}

\medskip\noindent
Now we can refine Question 3:

\medskip\noindent
{\bf Question 3.1.} Does Theorem
\ref{theorem-Bridgeland} describe all
geometric stability conditions?

\medskip\noindent
{\bf Question 3.2.} Do there exist
nongeometric stability conditions?

\medskip\noindent
{\bf Answer 3.2.}
\begin{itemize}
\item $\dim X=1$ no! Unless $X=\mathbb{P}^1$.
\item $\dim X \leq 3$? Yes if $X$ has a full
exceptional collection.
\item $\dim X=2$?
\begin{enumerate}
\item If $X$ is an abelian surface, no!
\item Yes for $\mathbb{P}^2$, and more
generally for rational surfaces.
\item Yes for K3 surfaces.
\item $X \supset C \cong \mathbb{P}^1$ such
that $C^2 < 0$.
\end{enumerate}
\end{itemize}

\medskip\noindent
{\bf Takeaway.} $\text{Stab}(X)$ sees
something about how (1) is different to
(2)-(4) but what is it?

\medskip\noindent
The following theorem is valid in any
dimension:

\begin{theorem}[Lie Fu-Chunyi Li-Xiaolei Zhao
'21]
\label{theorem-Albanese-variety}
$X$ has a finite Albanese morphism ($\iff$
there exist a finite morphism to an abelian
variety), then
$$
\text{Stab}(X) = \text{Stab}^{\text{geo}}(X)
$$
\end{theorem}

\begin{proof}[Sketch of proof]
\begin{itemize}
\item $\sigma \in \text{Stab}(X)$.
\item Let $E_1,\ldots,E_k$ be the {\it
Jordan-H\"older factors} of $\mathcal{O}_X$
w.r.t. $\sigma$ (i.e. $E_i$ $\sigma$-stable
connected component containing the identity).
\item $\mathcal{L} \in \text{Pic}^0(X)$.
\end{itemize}
By [Polishchuk '07], $\otimes \mathcal{L}$
does not change stability. Also
$\mathcal{O}_x \otimes \mathcal{L}\cong
\mathcal{O}_x$ so HN filtration and
Jordan-H\"older factors are preserved
$\implies
$ $E_i \otimes \mathcal{L} \cong E_i$.
\begin{itemize}
\item $\mathcal{L} \in
\text{Pic}^0(\text{Alb}(X))$, $E_i \otimes
\text{alb}^\vee(\mathcal{L})\cong E_i$.
\end{itemize}
$$
\begin{aligned}
\implies & R\text{alb}_*(E_i)\otimes
\mathcal{L} \cong R\text{alb}_*(E_i)\qquad
\text{(projection formula)}\\
\implies &  R\text{alb}_*(E_i)\text{ has
finite support}\qquad
\text{([Polishucuk '03])}\\
\implies &E_i\text{ has finite support}\qquad
\text{($\text{alb}_*$ finite)}\\
\overset{\text{claim}}{\implies
}&\mathcal{O}_x\text{ is $\sigma$-stable}
\end{aligned}
$$
\end{proof}

\medskip\noindent
{\bf Question 4.} Albanese morphism of $X$ is
not finite implies that there exist
nongeometric stability conditions?

\medskip\noindent
The idea is to investigate this for examples
arising as free quotients. If $G \mathbb{y}
X$ is a free action of a finite group, then
$\text{alb}_X$ is finite, and for $Y:=X/G$ we
have $\text{alb}_Y$ is not finite. In
sections 2 and 3 we will compate
$\text{Stab}(X)$ and $\text{Stab}(X/G)$.

\section{Equivariant categories}
\label{section-equivariant-categories}

Let $G$ be a finite group and $\mathcal{D}$ a
$k$-linear category (i.e. all $\Hom$ are
$k$-vector spaces, in particular it is
additive), where $k=\bar{k}$ and
$(|G|,\text{char}k)=1$.

The following definition is by [Deligne '97].

\begin{definition}[Group action]
\label{definition-group-action}
A {\it (right) action} of $G$ on
$\mathcal{D}$ is the data of
\begin{enumerate}
\item $\forall g \in G$,
$\phi_g:\mathcal{D}\xrightarrow{\sim}
\mathcal{D}$.
\item $\forall g,h\in G$, a natural
isomorphism
$\mathcal{E}_{g,h}:\phi_g\circ \phi_h
\xrightarrow{\sim}\phi_{hg}$ such that
$$
\xymatrix{
\phi_f \circ\phi_g \circ\phi_h
\ar[r]^{\mathcal{E}_{g,h}}\ar[d]_{\mathcal{E}
_{f,g}}&\phi_f
\circ \phi_{hg}
\ar[d]^{\mathcal{E}_{f \circ hg}}\\
\phi_{gf}\circ
\phi_h\ar[r]^{\mathcal{E}_{gf,h}}&\phi_{hgf}
}
$$
\end{enumerate}
\end{definition}

\begin{remark}
\label{remark-more-than-a-homomorphism}
This is more than a group a homomorphism $G
\to \text{Aut}(\mathcal{D})$. Indeed, given
this, $\forall E \in \mathcal{D}$,
$$
\phi_g \circ \phi_h(E)\cong \phi_{gh}(E)
$$
but may not come from a natural isomorphism
of functors
$\phi_g\circ\phi_h\cong\phi_{hg}$.
\end{remark}

\begin{example}
\label{example-group-actions}
\begin{enumerate}
\item $X$: scheme, $G \leq \text{Aut}(X)$.
For every $g \in G$ define
$\phi_g:=g^*
:\text{Coh}(X)\xrightarrow{\sim}\text{Coh}(X)
$.
Then for every $g,h$ there are canonical
isomorphisms
$$
\phi_g \circ\phi_h=g^*\circ
h^*\xrightarrow{\sim}(hg)^* =\phi_{hg}.
$$
$G \mathbb{y} \text{Coh}(X)$ lifts to
$G\mathbb{y}D^b(X)$.

Eg. an Enriques surface, i.e.
$Y=X/\mathbb{Z}/2\mathbb{Z}$, $X$ a K3
surface,
$\mathbb{Z}/2\mathbb{Z}=\left\langle{}i\right\rangle{}$, for
$i:X\to X$ involution, $x
\in X$, $i^*
(\mathcal{O}_x)=\mathcal{O}_{i^{-1}(x)}$.

\item $Q$: acyclic quiver, $G\leq
\text{Aut}^+(Q)=$ automorphisms of the
graph, orientation preserving. Then
$G\mathbb{y}\text{Rep}Q$. Via
$(M_i,\varphi_\alpha)\mapsto
(M_{g(i)},\varphi_{g(x)})$.

Eg. [Missing]
\end{enumerate}
\end{example}

\medskip\noindent
{\bf Equivariantization.}

\begin{definition}
\label{definition-G-equivariant-object}
Let $G \mathbb{y}\mathcal{D}$. A {\it
$G$-equivariant} object is a pair
$(E,\{\lambda_g\}_G$,
\begin{itemize}
\item $E \in \mathcal{D}$,
\item $\lambda_g:E
\xrightarrow{\sim}\phi_g(E)$, a {\bf choice}
of isomorphism
for all $g\in G$, such that
$$
\xymatrix{
E\ar[r]^{\lambda_g}\ar[d]_{\lambda_hg,\sim}&
\phi_g(E)\\
\phi_{hg}(E)&\phi_g(\phi_n(E))\ar[l]
_{\mathcal{E}_{gh}}
}
$$
A {\it morphism of $G$-equivariant objects}
is
$$
(E,\{\lambda_g\}_G)\to(E'_g,\{\lambda'_g\}_G)
$$
is $F \in \Hom_\mathcal{D}(E,E')$ such that
for all $g \in G$,
$$
\xymatrix{
E\ar[d]_F\ar[r]^{\lambda_g}&\phi_g(E)\ar[d]
^{\phi_g(F)}\\
E'\ar[r]^{\lambda'_g}&\phi_g(E')
}
$$
Together these form a category,
$\mathcal{D}_G$, called {\it $G$-equivariant
category}.
\end{itemize}
\end{definition}

\begin{remark}
\label{remark-uses-isomorphisms-Egh}
This uses the isomorphisms
$\mathcal{E}_{g,h}$.
\end{remark}

\begin{example}
\label{example-G-equivariant-objects}
\begin{enumerate}
\item $G\mathbb{y}X$, $G$-equivariant objects
are $G$-equivariant coherent
sheaves.
$$
(\text{Coh}(X))_G=:\text{Coh}_G(X)
$$
In fact, $\text{Coh}_G(X)\cong
\text{Coh}([X/G])$.

If $G$ acts freely, then
$$
\xymatrix{
X\ar[dr]_{\pi}\ar[rr]^g &  &
X\ar[dl]_{\pi}\\
&X/G
}
$$
and $\text{Coh}(X/G)\cong \text{Coh}_G(X)$
via $E \mapsto (\pi^* E,\{\lambda_g\}_{g \in
G}$ where
$$
\xymatrix{
\lambda_g:\pi^* E\ar[dr]_{\lambda_g}\ar[r]&
(\pi \circ g)^* E\ar[d]^=\\
&g^*(\pi^*E)
}
$$
Eg. $Y=X/(\mathbb{Z}/2\mathbb{Z})$ Enriques
surface, $y \in Y$,
$\text{supp}\pi^{-1}(y)=\{x,x'=i^{-1}(x)\}$
$$
\xymatrix@R=.2em{
&\mathcal{O}_x\ar[r]^{\text{id}}&\mathcal{O}
_x\\
\pi^*
(\mathcal{O}_x)=&\oplus&\oplus&=\text{id}
^*(\pi^*(\mathcal{O}_y))\\
&\mathcal{O}_{x'}\ar[r]^{\text{id}}&
\mathcal{O}_{x'}
}
$$
$$
\xymatrix@R=.2em{
&\mathcal{O}_x\ar[r]^{i^*}\ar[ddr]&
\mathcal{O}
_x\\
\pi^*
(\mathcal{O}_y)=&\oplus&\oplus&
=i^*(\pi^*(\mathcal{O}_y))\\
&\mathcal{O}_{x'}\ar[r]^{i^*}\ar[uur]&
\mathcal{O}_{x'}
}
$$
In general,
$$
(D^b(X))_G\cong
D^b(\text{Coh}_G(X))=:D^b_G(X)
$$
\item {\bf Theorem (Demonet '10).}
$G\mathbb{y} Q$ then $\exists Q_G$ such that
$\text{Rep}(Q_G)\cong\text{Rep}(Q)_G$. Eg.
[Missing.]

\item $G \mathbb{y} \mathcal{D}$ for $G$
abelian,
$\hat{G}:=\Hom(G,K^*)\mathbb{y}\mathcal{D}_G$
by
$$
\phi_{\underbrace{x}_{\in
G}}((E,\{\lambda_g\}_G))
:=(E,\{\lambda_g\}_G)\otimes
x=(E,\{\lambda_g\cdot x(g)\}_{g \in G})
$$
Eg. $Y$ Enriques surface \ldots{} [Missing]
\end{enumerate}
\end{example}

\begin{theorem}[Elagin '15]
\label{theorem-Elagin}
$G$ abelian, $G\mathbb{y}\mathcal{D}$
$$
(\mathcal{D}_G)_{\hat{G}}\cong \mathcal{D}
$$
\end{theorem}

\section{Exercise lecture}
\label{section-exercise-lecture}

On a surface $X$, to have a Bridgeland
stability condition we wanted a map
$$
\xymatrix{
Z:K_0(X)\ar[rr]\ar[dr]&  &  \mathbb{C}\\
&  \Lambda\ar[ur]
}
$$
For a numerical B.s.c., we ask
$$
\xymatrix{
Z:K_0(X)\ar[rr]\ar[dr]&  &  \mathbb{C}\\
&  K_{num}\ar[ur]
}
$$
Recall that the {\it Grothendieck group}
$K_0$ can be defined to be the free group
generated by isomorphism calsess of objects
in $D^b(X)$ quotient by the {\it Euler
pairing}, i.e.
$$
[F]=[E]+[G] \iff E \to F \to G \to E[i]
$$
Now let
$$
\chi(E,F)=\sum_{i}(-1)^{i}\dim
\underbrace{\text{Ext}^i(E,F)}
_{=\Hom(E,F[i])}
$$
then
$$
K_{num}=K_0(X)/\underbrace{\text{null}(\chi)}
_{=\{E:\chi(E,F)=0\forall F\in K_0(X)\}}
$$
To obtain a lattice consider
$$
\xymatrix{
K_0(X)\ar[d]\ar[r]^{ch}& \CH^* (X)\ar[d]\\
K_{num}(X)\ar[r]^{ch}&\underbrace{\CH^*_{num}
(X)}
_{\simeq \underbrace{\mathbb{Z}}_{ch_0}
\oplus \text{NS}(X) \oplus
\frac{1}{2}\underbrace{\mathbb{Z}}_{ch_2}}
}
$$
Fix $B \in \text{NS}(X)_{\mathbb{R}}$.
$$
\text{NS}(X)\ni A \mapsto \text{deg}(A\cdot
B)
\in
c\cdot \mathbb{Z}
$$
for $c \in \mathbb{R}$. For a surface,
$$
\Lambda = \mathbb{Z} \oplus
\mathbb{Z}^{\oplus \text{rk}(\text{NS}(X))}
\oplus \frac{1}{2}\mathbb{Z}
$$
$$
Z_{H,D,\alpha,\beta}=(\alpha+i \beta)H^2
ch_0+(D+iH)ch_1-ch_2.
$$

\begin{exercise}
\label{exercise-Ox-is-simple}
$\mathcal{O}_X$ is simple (i.e. no
subobjects) in $\text{Coh}^{H,\beta}(X)$.
\end{exercise}

\end{document}
